% Part: second-order-logic
% Chapter: syntax-and-semantics
% Section: expressive-power

\documentclass[../../../include/open-logic-section]{subfiles}

\begin{document}

\olfileid{sol}{syn}{exp}
\olsection{اظہار دی طاقت}

\begin{explain}
دوجے درجے دے متغیراں اُتے مسوّریت زبان دی
اظہار دی طاقت نوں پہلے درجے دی منطق دے مقابل
بہت ودھا دیندی اے۔ دوجے درجے دی وجودی مسوّریت
نال اسیں آکھ سکدے آں کہ کجھ خاصیتاں والے فنکشن یا تعلق موجود نیں۔ پہلے درجے دی منطق
وچ ایہہ کہن دا اکلوتا طریقہ ایہہ اے کہ اوس مقصد
لئی اک غیر منطقی علامت، یعنی !!a{function}
یا !!{predicate}، مقرر کیتی جائے۔ دوجے درجے
دی عمومی مسوّریت نال اسیں آکھ سکدے آں کہ
دائرے دے سارے ذیلی سیٹ، اُس اُتے سارے
تعلق، یا !!{domain} توں !!{domain} ول سارے
فنکشوناں وچ کوئی خاصیت اے۔ پہلے درجے دی
منطق وچ اسیں صرف اوہناں ذیلی سیٹاں، تعلقاں
یا فنکشوناں بارے ایہہ آکھ سکدے آں جیہڑے زبان
دی کسے غیر منطقی علامت نوں دتے گئے ہون۔
تے جدوں اسیں کسے خاصیت والے ذیلی سیٹاں،
تعلقاں یا فنکشوناں دے وجود، یا انہاں سبھناں
دے اوہ خاصیت رکھن، بارے آکھدے آں، تاں اُس
خاصیت نوں بیان کردیاں وی دوجے درجے دی مسوّریت
ورت سکدے آں۔ ایس نال اسیں اوہناں تعلقاں دی تعریف کر سکدے آں
جیہڑے پہلے درجے دی منطق وچ قابلِ تعریف نہیں، تے دائرے دیاں اوہ خاصیتاں
بیان کر سکدے آں جو پہلے درجے دی منطق وچ
بیان نہیں ہو سکدیاں۔
\end{explain}

\begin{defn}
جے $\Struct{M}$ زبان~$\Lang{L}$ لئی
!!a{structure} اے، تاں تعلق~$R \subseteq
\Domain{M}^2$،~$\Lang{L}$ وچ
\emph{قابلِ تعریف} اے جے کوئی !!{formula}~$!A_R(x, y)$
ایہو جیہا ہووے جیہدے وچ صرف $x$
تے~$y$ آزاد ہون، تے $R(a, b)$
(یعنی $\tuple{a,b} \in R$) اُدوں تے صرف
اُدوں لاگو اے جدوں $\Sat{M}{!A_R(x, y)}[s]$
اوہ قدر بندی لئی جیہڑی $s(x) = a$
تے $s(y) = b$ رکھے۔
\end{defn}

\begin{ex}
پہلے درجے دی منطق وچ اسیں برابری دا
تعلق~$\Id{\Domain{M}}$ (یعنی
$\Setabs{\tuple{a,a}}{a \in
  \Domain{M}}$) فارمولے~$\eq[x][y]$
نال متعین کر سکدے آں۔ دوجے درجے دی منطق وچ
ایہہ تعلق \emph{$\eq$ توں بناں} وی متعین ہو سکدا اے۔ جے $a$ تے $b$،~$\Domain{M}$
دے اکّو !!{element} ہون، تاں اوہ
~$\Domain{M}$ دے اوہناں ہی ذیلی سیٹاں دے
!!{element}s ہون گے (کیوں جے سیٹ اپنے
!!{element}s نال متعین ہوندے نیں)۔
اُلٹے پاسے، جے $a$ تے $b$ وکھرے ہون،
تاں کوئی ذیلی سیٹ ایسا اے جیہدے !!{element}s وچ اک اے
تے دوجا نہیں۔ مثال لئی $a \in \{a\}$،
پر $b
\notin \{a\}$ جے $a \neq b$۔
سو «~$\Domain{M}$ دے اکّو ذیلی سیٹاں دے
!!{element}s ہونا» اوہ تعلق اے جیہڑا
$a$ تے $b$ دے لئی اُدوں تے صرف اُدوں
درست اے جدوں $a = b$۔ ایہہ تعلق
دوجے درجے دی منطق وچ بیان ہو سکدا اے،
کیوں جے اسیں~$\Domain{M}$ دے سارے
ذیلی سیٹاں اُتے مسوّریت کر سکدے آں۔
سو ہِٹھلا !!{formula}،
$\Id{\Domain{M}}$ نوں متعین کردا اے:
\[
\lforall[X][(X(x) \liff X(y))]
\]
\end{ex}

\begin{prob}
دکھاؤ کہ $\lforall[X][(X(x) \lif X(y))]$
(دھیان: $\lif$ اے،~$\liff$ نہیں!)
$\Id{\Domain{M}}$ نوں متعین کردا اے۔
\end{prob}

\begin{ex}
جے $R$ اک دو رُکنی !!{predicate} اے، تاں
$\Assign{R}{M}$،~$\Domain{M}$ اُتے
اک دو رُکنی تعلق اے۔ شاید کجھ الجھن
پیدا ہووے، پر اسیں $R$ نوں !!{predicate}~$R$
تے خود تعلق~$\Assign{R}{M}$ دوناں لئی
ورتاں گے۔ \emph{تعدی تکمیل}~$R^*$، جیہڑی~$R$
دی اے، اوہ تعلق اے جو $a$ تے $b$ دے درمیان
اُدوں تے صرف اُدوں ہوندا اے جدوں کجھ
$c_1$, \dots, $c_k$ لئی
$R(a,c_1)$, $R(c_1, c_2)$, \dots,
$R(c_k,b)$ درست ہون۔ ایس وچ $k = 0$
والی صورت وی شامل اے؛ یعنی جے
$R(a,b)$ ہوندا اے، تاں $R^*(a,b)$
وی ہوندا اے۔ ایس دا مطلب
$R \subseteq R^*$ اے۔ اصل وچ
$R^*$ اوہ سب توں چھوٹا متعدی تعلق اے
جیہڑا~$R$ نوں شامل کردا اے۔
دوجے درجے دی منطق وچ
اسیں ایہہ آکھ سکدے آں کہ $X$ اک
متعدی تعلق اے جیہڑا~$R$ نوں شامل
کردا اے:
\begin{multline*}
  !B_R(X) \ident \lforall[x][\lforall[y][(R(x,y) \lif X(x, y))]] \land {}\\
\lforall[x][\lforall[y][\lforall[z][((X(x,y) \land X(y,z)) \lif X(x,
      z))]]].
\end{multline*}
پہلا عطف آکھدا اے کہ $R \subseteq X$،
تے دوجا آکھدا اے کہ $X$ متعدی اے۔

ایہہ آکھنا کہ $X$ ایہو جیہا سب توں
چھوٹا تعلق اے، اِس گل دے برابر اے کہ
اوہ خود ہر اوس تعلق وچ شامل اے جیہڑا
$R$ نوں شامل کردا تے متعدی اے۔
سو اسیں~$R$ دی تعدی تکمیل ہیٹھلے
!!{formula} نال متعین کر سکدے آں:
\[
R^*(X) \ident !B_R(X) \land \lforall[Y][(!B_R(Y) \lif
  \lforall[x][\lforall[y][(X(x, y) \lif Y(x,y))]])].
\]
$\Sat{M}{R^*(X)}[s]$ اُدوں تے صرف اُدوں لاگو اے
جدوں $s(X) = R^*$۔
~$R$ دی تعدی تکمیل پہلے درجے دی
منطق وچ بیان نہیں ہو سکدی۔
\end{ex}

\end{document}
